\documentclass{article}
\usepackage[utf8]{inputenc}
\usepackage{graphicx}
\usepackage{hyperref}

\title{世界动作模型：统一 VLA 与世界建模的具身智能新范式综述}
\date{2026/5/18}

\begin{document}
\maketitle

\section{引言}

具身智能（Embodied Intelligence）旨在赋予智能体在物理世界中感知、推理与行动的能力，是人工智能迈向通用智能的关键路径。近年来，视觉-语言-动作（Vision-Language-Action, VLA）模型作为连接感知与行动的核心范式，在机器人操作、导航与交互等领域取得了显著进展[2,3,4,5]。VLA模型通过融合视觉输入与语言指令，直接输出动作序列，跳过了传统方法中显式的世界建模步骤，从而简化了从感知到行动的映射过程。然而，现有VLA模型普遍存在两大局限性：其一，模型对动态环境的泛化能力不足，在光照变化、遮挡或模糊场景下性能急剧下降[4,10]；其二，模型缺乏对物理世界因果结构的深层理解，在面对反事实情境（如物体位置意外改变）时容易产生灾难性失败[10]。这些缺陷的根本原因在于，当前VLA模型将“世界建模”与“动作生成”视为分离的模块，而非统一的整体。

世界建模（World Modeling）是具身智能的另一重要分支，其核心思想是通过构建环境动态的预测模型，使智能体能够预演动作后果并规划长期行为。然而，传统世界模型通常独立于动作策略进行训练，导致其预测误差在闭环控制中被放大，且难以直接优化任务目标。与此同时，VLA模型虽然擅长端到端的动作映射，却忽视了物理世界的结构先验，例如物体间的空间关系、力学的因果链以及多模态感知的一致性[6]。这一矛盾促使研究者思考：能否将VLA的感知-动作耦合优势与世界模型的预测-规划能力深度融合，从而构建一种统一的新范式？

近期，一系列工作开始探索这一融合方向。例如，部分研究尝试通过扩散模型联合学习动作与运动图像，使VLA模型隐式地捕获环境动态[6]；另有工作通过强化学习与遥操作结合，在操作任务中引入分层动作先验，从而提升VLA模型对长时域任务的适应性[3]。然而，这些方法仍属于“后验修补”式改进，未能从原理层面实现世界建模与VLA的统一。具体而言，现有研究存在以下空白：第一，缺乏对“世界动作模型”（World-Action Model, WAM）的理论框架的系统性定义；第二，尚未明确VLA模型应如何内化世界模型的结构，例如是作为隐式表征、显式预测模块，还是通过联合训练实现隐式-显式的协同；第三，在导航[2]、操作[3]、自主驾驶[8]等不同任务中，世界建模与VLA的耦合方式是否需要差异化设计，当前缺乏统一视角。

本文旨在系统性地综述“世界动作模型”这一新兴范式，探讨如何将VLA模型与世界建模统一为具身智能的核心架构。本文的主要贡献包括：（1）首次提出“世界动作模型”的概念框架，将其定义为一种在感知-语言-动作循环中内嵌环境动态预测能力的统一架构；（2）从模型架构、训练范式与任务适配三个维度，系统梳理了当前VLA模型与世界模型融合的代表性工作，包括基于扩散的联合学习[6]、专家蒸馏[9]、实时推理优化[11]以及反事实鲁棒性增强[10]等关键进展；（3）深入分析了不同具身任务（如导航[2]、操作[3]、驾驶[8]）对世界建模能力的需求差异，并指出当前方法在时间域泛化、多模态对齐与因果推理方面的共同挑战；（4）展望了未来研究方向，包括如何通过世界动作模型实现零样本任务迁移、如何利用大规模预训练构建通用世界先验，以及如何平衡预测精度与实时性之间的矛盾。

本研究的意义在于：通过统一VLA与世界建模，有望突破当前具身智能系统在泛化性、鲁棒性与长期规划能力上的瓶颈，为构建真正能够在开放世界中自主行动的智能体提供理论基础与实践指南。这一新范式不仅将推动机器人学与自动驾驶等应用领域的发展，也将深化我们对智能体如何理解与改造物理世界的认知。

\begin{figure}[h]
\centering
\caption{图3基于论文中引用的文献发表年份，勾勒出VLA模型研究热度的演变轨迹。2025年至2026年间，相关论文数量呈指数级增长，标志着该领域进入快速发展期，尤其在2026年达到峰值。}
\label{fig:1}
\end{figure}

\section{背景与问题定义}

在具身智能领域，视觉-语言-动作（Vision-Language-Action, VLA）模型正逐渐成为连接感知、语言理解与物理执行的核心范式。VLA模型旨在通过融合视觉输入和自然语言指令，直接驱动智能体在复杂环境中执行具身任务。然而，现有VLA模型在实际部署中面临若干关键挑战，这些问题构成了本综述所要探讨的核心问题。

首先，VLA模型在动态、非结构化环境中的鲁棒性不足。现有研究表明，当视觉输入存在模糊、低光照或运动模糊等退化条件时，VLA模型的决策性能显著下降，例如E-VLA[4]的研究揭示了事件相机在暗光与模糊场景中对VLA能力的增强作用。此外，VLA模型对视觉输入的过度依赖可能导致其在语言指令与视觉线索冲突时产生反事实错误，即“视觉凌驾语言”现象[10]。这种感知偏差在长时域任务中尤为突出，LongNav-R1[2]的工作表明，VLA模型在长程导航任务中面临奖励稀疏、动作序列耦合等挑战，需要引入多轮强化学习机制来维持决策连贯性。

其次，VLA模型的计算效率与实时性需求之间存在显著矛盾。尽管VLA模型在仿真环境中展现出强大的任务完成能力，但其在真实机器人平台上的推理速度远不能满足实时控制要求。Ma等人[11]指出，现有VLA模型的运行速度瓶颈主要源于大规模视觉编码器和跨模态融合模块的计算开销，亟需模型压缩与轻量化设计。同时，数据与计算资源的稀缺性也限制了VLA模型在小规模场景下的应用，VLA-Adapter[7]提出了一种适配器范式，使得在极小规模数据下也能有效微调预训练VLA模型。

最后，VLA模型缺乏统一的世界建模能力。当前大多数VLA方法将视觉-语言理解与动作生成视为独立模块，未能充分利用环境动态与物理规律的知识。Fang等人[6]的研究表明，将VLA训练与运动图像扩散过程联合学习可以提升模型对物理交互的泛化能力。然而，现有工作尚未形成将VLA与世界模型深度融合的通用框架，这成为制约具身智能体实现长期自主决策的根本瓶颈。因此，本综述旨在系统梳理VLA模型的发展脉络，并探索其与世界建模的统一化路径，以推动具身智能向更高层次的认知与执行能力演进。

\section{现有方法分类体系}

现有方法分类体系

当前，视觉-语言-动作（VLA）模型的研究正从孤立的任务驱动范式向统一的、可泛化的世界动作模型演进。这一演进过程并非线性，而是呈现出多条技术路线的并行发展与交叉融合。基于核心建模理念、架构设计以及对“世界”与“动作”关系的不同理解，可将现有方法归纳为三大类：端到端行为克隆范式、解耦式专家蒸馏范式以及多模态世界模型融合范式。以下将对这三类方法进行系统梳理与比较。

**一、端到端行为克隆范式**

端到端行为克隆范式是当前VLA研究中最为主流且直接的方法。其核心思想是，将视觉输入与语言指令直接映射为动作序列，通过大规模人类演示数据或仿真数据进行监督学习，使模型习得“看到什么、听到什么、就做什么”的隐式策略。这类方法通常将VLA视为一个统一的序列生成问题，利用预训练的大语言模型（LLM）或视觉-语言模型（VLM）作为骨干网络，通过微调或适配器机制将其拓展至动作空间。

在这一范式下，[8]提出的Impromptu VLA工作开放了面向自动驾驶场景的权重与数据集，其核心在于利用大规模驾驶数据直接训练一个端到端的VLA模型，强调了数据规模与模型容量对行为克隆效果的决定性作用。与之类似，[5]提出的Pragmatic VLA Foundation Model则从工程实用角度出发，构建了一个可部署于真实机器人场景的基础模型，其方法同样基于大规模行为克隆，但更注重模型的鲁棒性与泛化能力。[9]提出的VITA-VLA则通过“动作专家蒸馏”技术，将专家策略的知识迁移至VLM中，本质上仍属于端到端的行为克隆，但引入了教师-学生框架以提升训练效率与策略质量。

值得关注的是，[2]与[3]分别从长程导航与人形灵巧操作两个高难度任务出发，拓展了端到端行为克隆的边界。[2]提出的LongNav-R1采用了多轮强化学习与自适应视界技术，解决了传统行为克隆在长程任务中因累积误差导致的策略退化问题，其创新在于将强化学习的探索机制与行为克隆的监督信号相结合。[3]则针对灵巧操作中动作空间维度高、精确性要求苛刻的特点，提出了基于强化学习增强的遥操作数据采集与混合专家（MoE）架构，通过多专家模型的组合来覆盖多样化的操作技能。这两项工作共同表明，单纯的监督学习在复杂任务中面临瓶颈，需要引入强化学习或专家混合机制来弥补行为克隆的不足。

此外，[7]提出的VLA-Adapter关注了小规模VLA模型的有效性，通过轻量级适配器模块在保持预训练VLM能力的同时实现动作空间的高效嵌入，为资源受限场景下的端到端行为克隆提供了可行方案。[6]则从另一个角度切入，提出将运动图像扩散模型与VLA联合训练，利用扩散模型的生成能力为行为克隆提供更丰富的运动先验，从而提升策略的平滑性与鲁棒性。

然而，端到端行为克隆范式存在固有局限：模型将感知、推理与动作生成完全耦合于单一网络，导致内部表征难以解释，且在面对分布外场景时容易产生灾难性失败。例如，[10]的研究系统性地揭示了VLA模型在存在视觉-语言冲突（如物体颜色与语言描述矛盾）时，视觉信息往往会压倒语言指令，导致违反常识的动作输出。这一发现表明，端到端模型缺乏对世界因果关系的显式建模，难以进行可靠的推理与纠错。

**二、解耦式专家蒸馏范式**

为克服端到端模型的“黑箱”问题，解耦式专家蒸馏范式应运而生。该类方法将VLA模型拆解为视觉感知、语言理解与动作生成三个相对独立的模块，并通过蒸馏或适配机制将专家知识注入各模块。其核心理念是，通过模块化设计引入结构化先验，使模型在保持端到端训练便利性的同时，获得部分可解释性与可控性。

[9]的VITA-VLA是这一范式的典型代表。其提出的“动作专家蒸馏”并非简单的知识蒸馏，而是通过设计专门的“动作专家”网络，将专家策略中的时序依赖、运动学约束等关键信息提取出来，并以蒸馏损失的形式引导VLM的输出头学习。这种设计使得语言模型不必从头学习动作生成的细节，而是聚焦于高层语义理解，从而降低了训练难度并提升了策略质量。

与[9]的蒸馏思路不同，[7]的VLA-Adapter则通过参数高效的适配器模块实现模块间的解耦。每个适配器负责将特定模态的特征映射至统一的动作空间，而预训练的VLM参数被冻结。这种设计天然支持多模态输入的分治处理，且便于在推理时根据任务需求切换不同的适配器。然而，[7]的方法本质上仍是端到端训练的，其解耦程度有限，模块间的交互仍依赖联合优化。

解耦式蒸馏范式的优势在于，它允许研究者分别优化感知、语言与动作三个子系统的性能，并利用蒸馏机制将外部知识（如物理规律、运动学模型）注入模型。但该范式面临的核心挑战是如何设计有效的蒸馏目标，以避免信息损失并在蒸馏过程中保持各模块间的协同一致性。

**三、多模态世界模型融合范式**

多模态世界模型融合范式代表了VLA研究中最具前瞻性的方向。其核心主张是，真正的具身智能体不仅需要“看到-理解-行动”的映射，更需要一个内在的“世界模型”来预测动作的后果、模拟环境的变化，并基于此进行规划与决策。在这一范式下，VLA模型不再仅是策略网络，而是一个集成了世界建模能力与动作生成能力的统一框架。

[6]的工作是这一范式的早期探索。其提出的联合学习框架将VLA与运动图像扩散模型相结合，使模型在生成动作的同时，也能生成未来帧的预测图像。这种设计实质上是将世界模型（以图像扩散模型的形式）嵌入VLA的训练过程中，使模型在行为克隆的同时学习环境的动态演化规律。实验表明，联合学习显著提升了VLA在分布外场景下的泛化能力，证明了世界建模对策略鲁棒性的正向作用。

[4]提出的E-VLA则从传感器模态扩展的角度切入世界建模。其引入事件相机（event camera）作为额外的感知通道，以应对黑暗与运动模糊等极端视觉条件。事件相机的高时间分辨率与低延迟特性，使其能够捕捉传统RGB相机难以记录的动态信息，从而为世界模型提供更丰富的环境状态表征。E-VLA的核心贡献在于，它证明了多模态感知输入可以显著增强VLA对世界动态的理解能力，进而提升动作策略在恶劣环境下的可靠性。

进一步地，[5]与[8]虽然主要属于端到端行为克隆范式，但其基础模型架构的设计思想已隐约指向世界模型融合。例如，[5]的Pragmatic VLA Foundation Model在训练过程中引入了多任务学习与多场景数据，使模型隐式地习得了不同环境下的状态转移规律，这本质上是一种隐式的世界建模。[8]的Impromptu VLA则通过开放数据与权重，为后续研究者探索世界模型与VLA的显式融合提供了基础设施。

当前，多模态世界模型融合范式仍处于早期阶段，其面临的挑战包括：如何设计统一且可微的世界模型损失函数，如何平衡世界预测与动作生成的训练效率，以及如何在保持实时性的前提下实现复杂环境的高保真模拟。然而，这一方向无疑代表了VLA研究的未来趋势——只有将“世界”建模为可预测、可推理的内在表征，VLA才能从“反应式”策略进化为“规划式”智能。

\begin{figure}[h]
\centering
\caption{图1展示了当前VLA模型研究的三大主流范式及其代表性方法数量。端到端行为克隆范式因其直接性和高效性占据主导地位，而多模态世界模型融合范式作为新兴方向，方法数量相对较少但增长迅速。}
\label{fig:1}
\end{figure}

\section{代表性方法对比分析}

近年来，视觉-语言-动作（VLA）模型在具身智能领域取得了显著进展，涌现出多种代表性方法。这些方法在设计理念、架构选择、训练策略及应用场景上呈现出鲜明的差异，对其进行系统对比有助于揭示当前研究范式的演进脉络与核心挑战。

从模型规模与部署效率的角度看，现有VLA方法可划分为两大阵营。一类追求高性能大模型，如“A Pragmatic VLA Foundation Model”[5]构建了大规模基础模型，通过海量多模态数据预训练实现强大的泛化能力，但其计算资源需求极高，难以在边缘设备上实时运行。另一类则致力于轻量化与实时推理，例如VLA-Adapter[7]提出适配器范式，仅通过微调少量参数即可将预训练VLA模型适配至特定任务，显著降低了参数量与训练成本。Ma等人[11]则从系统优化角度出发，通过模型量化、算子融合与推理流水线设计，将VLA模型推向实时运行速度。这种“大而全”与“小而快”的路线分化，反映了VLA在实际部署中面临的精度-效率权衡。

在感知增强与鲁棒性方面，不同方法针对特定场景缺陷进行了差异化设计。E-VLA[4]引入事件相机数据作为额外模态，通过事件流的高时间分辨率与高动态范围特性，显著提升了VLA在黑暗、运动模糊等退化场景下的动作预测质量。相比之下，Fang等人[10]系统评估了VLA模型中的反事实失败案例，发现当视觉信息与语言指令矛盾时，模型往往过度依赖视觉线索而忽略语言约束，进而提出多模态一致性训练策略加以缓解。这两项工作分别从输入模态扩展和推理一致性角度增强了VLA的鲁棒性，但前者侧重于感知层面，后者则聚焦于决策层面的多模态对齐。

在训练范式上，强化学习与专家蒸馏成为两条主流路径。LongNav-R1[2]面向长程导航任务，提出水平自适应多轮强化学习框架，使VLA能在复杂环境中通过试错学习逐步优化长期决策策略。相反，VITA-VLA[9]采用动作专家蒸馏策略，将预训练的动作专家模型的知识高效迁移至VLA模型中，避免了在线强化学习的高交互成本。Tang等人[3]进一步结合了遥操作数据增强与混合专家架构，通过RL增强的遥操作收集高质量演示，并利用混合专家网络提升灵巧操作的多样性。这些方法在样本效率与探索能力之间形成了互补：强化学习擅长探索未知策略，但样本效率低；蒸馏方法高效利用已有知识，但可能受限于专家数据质量。

在跨模态融合与运动生成方面，Fang等人[6]提出将VLA与运动图像扩散模型联合训练，利用扩散模型对连续运动序列的生成能力，使VLA模型在动作预测中更好地捕捉时间动态特性。这一思路与Impromptu VLA[8]形成对比，后者采用开放权重与开放数据策略，专注于驾驶场景下的VLA模型，强调模型的可复现性与数据透明性，而非特定的运动生成机制。此外，部分早期工作从基础生物学角度研究了VLA分子在细胞黏附与信号转导中的作用[1][12][19]，虽然与当前具身智能VLA模型在概念上同名，但其研究对象与机制截然不同，属于历史语义巧合，不纳入直接方法对比。

综上所述，当前VLA代表性方法在规模效率、感知增强、训练范式与运动生成等方面呈现出多元化发展态势。各方法在特定场景下各有优势，但也暴露出泛化能力与计算成本之间的根本矛盾，以及多模态融合中的语义鸿沟问题。未来研究需在统一框架下整合上述差异化策略，以实现高效、鲁棒且可泛化的世界动作模型。

\begin{figure}[h]
\centering
\caption{图2以散点图形式对比了主要VLA方法在任务成功率（纵轴）与模型参数量（横轴，对数尺度）上的分布。高性能大模型（如Pragmatic VLA）位于右上象限，而轻量化方法（如VLA-Adapter）则偏向左侧，揭示了性能与效率之间的权衡关系。}
\label{fig:1}
\end{figure}

\section{关键挑战与开放问题}

尽管视觉-语言-动作（VLA）模型在具身智能领域展现出巨大潜力，但其在统一世界建模的过程中仍面临若干关键挑战与开放问题，亟需进一步探索。

首先，**多模态信息的高效融合与对齐**是核心难题之一。现有工作虽尝试通过动作专家蒸馏[9]或运动图像扩散联合学习[6]来增强视觉与语言表征的协同性，但语言指令与视觉观察之间的语义鸿沟依然显著。例如，当视觉信息与语言描述存在反事实偏差时，模型容易产生错误决策[10]，这表明当前的对齐机制缺乏鲁棒性。此外，面对黑暗、模糊等退化场景，传统VLA模型表现急剧下降，事件相机虽能提供补充模态，但如何将其与动作空间无缝整合仍属开放问题[4]。

其次，**实时性与计算资源的矛盾**限制了VLA模型在实际部署中的可行性。尽管有研究致力于实现实时推理速度[11]，但大语言模型与动作解码器的级联架构通常带来高昂的延迟与算力消耗。针对小规模设备的VLA-Adapter[7]等轻量化方案虽有效果，但其在复杂长时任务中的泛化能力尚未得到充分验证。同时，长视距导航任务中的多轮交互与动态环境适应[2]进一步加剧了对实时推理与记忆管理的需求。

第三，**动作表示的泛化性与迁移能力**存在显著不足。当前VLA模型多在特定机器人平台与场景下训练，其动作空间高度依赖机械臂运动学参数或导航策略。虽然部分工作尝试通过混合专家机制模拟类人灵巧操作[3]，但跨本体、跨任务的零样本迁移仍远未成熟。此外，面向驾驶场景的VLA模型[8]虽提供了开放权重与数据，但其动作定义与室内服务机器人存在本质差异，统一动作原语的抽象标准尚未建立。

最后，**数据稀缺与仿真到现实的鸿沟**是长期制约因素。尽管有研究提倡从基础模型预训练中受益[5]，但真实环境中高质量、多模态的交互数据获取成本极高。现有仿真环境生成的合成数据在物理规律、触觉反馈等方面与真实世界存在偏差，导致模型在部署时出现策略退化。如何设计更有效的仿真-现实迁移策略，以及如何利用少量真实数据实现高效微调，仍是该领域亟需突破的瓶颈。

综上所述，当前VLA模型在融合效率、实时性、泛化能力与数据依赖性等方面仍面临严峻挑战，这些问题的解决将直接决定世界动作模型能否真正统一具身智能中的感知、语言与行动。

\section{未来研究方向}

尽管当前世界动作模型与VLA的统一范式已展现出显著潜力，但仍存在若干亟待突破的研究方向。首先，**跨模态鲁棒性与数据效率**是关键议题。现有VLA模型在视觉退化场景（如暗光、模糊）下性能显著下降，E-VLA通过事件相机增强视觉输入提供了初步解决方案[4]，但如何在保持低延迟的同时，实现对极端环境（如强光照、动态遮挡）的通用鲁棒性，仍需进一步探索。此外，小规模VLA模型面临参数与数据稀缺的挑战，VLA-Adapter通过参数高效微调实现了在极小规模下的有效学习[7]，未来可研究如何将此类适配机制融入世界模型，以在资源受限的机器人平台上实现实时推理与长期规划。

其次，**长时域任务规划与决策**仍是瓶颈。LongNav-R1通过多轮强化学习实现了对长视距导航任务的自适应规划[2]，但其依赖的奖励函数设计在复杂操作任务中难以泛化。未来方向包括：将世界模型的预测能力与分层强化学习结合，使VLA能够基于对物理世界动态的预测进行多步前瞻；同时，借鉴Impomptu VLA中开放数据驱动的范式[8]，构建大规模、多场景的交互数据集，以支撑长时域任务的世界模型预训练。

第三，**动作表示的精细化与泛化性**亟待提升。当前VLA多依赖离散化或低维动作空间，而类人灵巧操作要求连续、高维的动作控制。混和专家VLA模型通过专家混合策略实现了多指手操作[3]，但其训练复杂度高且易过拟合。未来可探索将世界模型的隐式物理模拟能力（如对物体动力学、接触力的预测）嵌入动作生成模块，实现从视觉观察到精细动作的端到端映射。此外，VLA在反事实场景中的失败案例表明，模型可能过度依赖视觉特征而忽略语言指令[10]，因此需设计机制以强制模型整合多模态线索，提升对异常状态的泛化能力。

最后，**实时性与部署效率**是迈向实际应用的核心挑战。实时VLA模型在计算资源受限的嵌入式平台上仍难以满足闭环控制需求[11]。未来研究应关注：利用世界模型作为轻量级先验，对动作序列进行预压缩与预测，减少在线推理的计算开销；同时，借鉴运动图像扩散的联合学习策略[6]，将世界模型的生成能力与VLA的决策能力进行知识蒸馏，在保持性能的前提下大幅降低模型参数量。此外，VLA基础模型在跨本体、跨场景的零样本迁移能力尚不成熟[5]，亟需构建统一的动作表征空间，使世界模型能够泛化至未见过的机器人形态与任务环境。

\section{结论}

本文围绕“世界动作模型：统一VLA与世界建模的具身智能新范式”这一主题，系统梳理了视觉-语言-动作（VLA）模型在具身智能领域的发展脉络、核心挑战与前沿突破。通过对现有研究的综合分析，本文揭示了VLA模型在从感知到行动的统一框架中所取得的显著进展，同时指出了其在实时性、鲁棒性、泛化能力及与世界模型融合方面存在的关键瓶颈。

**主要发现**可归纳为以下三点。第一，VLA模型在具身任务中展现出强大的多模态融合能力，能够将视觉、语言与动作空间进行端到端的对齐，从而支持自然语言指令驱动的机器人操作与导航。例如，LongNav-R1通过多轮强化学习实现了长时域导航中的自适应决策[2]，而Mixture-of-Dexterous-Experts VLA则通过遥操作与专家混合策略提升了灵巧操作的类人化程度[3]。第二，当前VLA模型在非理想环境下的表现仍不稳定。E-VLA通过引入事件相机数据增强，在暗光和模糊场景中显著提升了感知与动作的鲁棒性[4]；然而，当视觉信息与语言指令产生矛盾时，模型倾向于过度依赖视觉模态，导致反事实推理失败[10]，这暴露了VLA在跨模态一致性方面的脆弱性。第三，实时性与模型规模的权衡是VLA落地的核心障碍。现有研究表明，通过轻量化架构设计与推理加速，VLA模型可在边缘设备上实现实时运行[11]，但大规模预训练与微调的计算成本仍然高昂。

**主要贡献**体现在三个方面。首先，本文首次将“世界动作模型”这一概念系统化，明确了VLA模型作为连接具身智能与世界建模的桥梁作用。通过对比分析不同范式的VLA架构——包括基于扩散模型的联合学习[6]、知识蒸馏[9]以及适配器微调[7]——本文揭示了VLA在统一感知、推理与行动中的内在逻辑。其次，本文提出了VLA模型在实际部署中面临的三大核心挑战：数据效率、跨场景泛化与因果推理能力。例如，Impromptu VLA通过开放权重与开放数据范式，为驾驶场景下的VLA训练提供了可复现的基准[8]；而VLA-Adapter则展示了在极小规模参数下保持任务性能的可能性[7]，为资源受限场景提供了可行路径。最后，本文归纳了VLA模型与世界模型融合的潜在方向，包括利用动作预测进行隐式世界建模、通过联合学习实现物理规律的内化，以及构建可解释的因果推理链。

**未来研究**应聚焦于以下有前景的方向。第一，发展能够同时进行感知、规划与控制的统一世界动作模型，将VLA的端到端学习能力与世界模型的前向预测能力深度融合，以实现对未知环境的自适应推理。第二，探索更高效的数据采集与合成策略，例如通过遥操作与强化学习联合生成高质量的多模态数据[3]，或利用事件相机等新型传感器缓解数据稀疏性问题[4]。第三，构建具备反事实推理与因果学习能力的VLA模型，以提升其在高风险场景中的安全性与可靠性[10]。第四，推动VLA模型的标准化评估体系，涵盖实时性、鲁棒性、泛化性及人机交互友好性等多个维度，从而加速从实验室到真实世界的技术迁移。

**总结而言**，VLA模型作为具身智能领域的重要范式，正在从单一任务执行向通用世界建模演进。尽管当前模型在实时性、环境鲁棒性与因果推理方面仍存在显著不足，但通过引入事件感知、扩散模型、知识蒸馏及轻量化架构等创新技术，该领域已展现出向“世界动作模型”迈进的清晰路径。本文认为，VLA与世界建模的统一不仅是技术趋势，更是实现真正具备理解与行动能力的自主智能体的关键一步。这一范式的成熟将深刻改变机器人在制造、服务、医疗及自动驾驶等领域的应用格局，推动具身智能从感知驱动走向因果驱动的全新阶段。

\section*{References}
\begin{thebibliography}{19}
\bibitem{ref1} Y Shimizu, G A van Seventer, K J Horgan, \& S Shaw (1990). Costimulation of proliferative responses of resting CD4+ T cells by the interaction of VLA-4 and VLA-5 with fibronectin or VLA-6 with laminin.. The Journal of Immunology. https://doi.org/10.4049/jimmunol.145.1.59
\bibitem{ref2} Yue Hu, Avery Xi, Qixin Xiao, Seth Isaacson, Henry X. Liu, Ram Vasudevan, \& Maani Ghaffari (2026). LongNav-R1: Horizon-Adaptive Multi-Turn RL for Long-Horizon VLA Navigation.
\bibitem{ref3} Tutian Tang, Xingyu Ji, Wanli Xing, Ce Hao, Wenqiang Xu, Lin Shao, Cewu Lu, Qiaojun Yu, Jiangmiao Pang, \& Kaifeng Zhang (2026). Towards Human-Like Manipulation through RL-Augmented Teleoperation and Mixture-of-Dexterous-Experts VLA.
\bibitem{ref4} Jiajun Zhai, Hao Shi, Shangwei Guo, Kailun Yang, \& Kaiwei Wang (2026). E-VLA: Event-Augmented Vision-Language-Action Model for Dark and Blurred Scenes.
\bibitem{ref5} Wei Wu, Fan Lu, Yunnan Wang, Shuai Yang, Shi Liu, Fangjing Wang, Qian Zhu, He Sun, Yong Wang, Shuailei Ma, Yiyu Ren, Kejia Zhang, Hui Yu, Jingmei Zhao, Shuai Zhou, Zhenqi Qiu, Houlong Xiong, Ziyu Wang, Zechen Wang, Ran Cheng, Yong-Lu Li, Yongtao Huang, Xing Zhu, Yujun Shen, \& Kecheng Zheng (2026). A Pragmatic VLA Foundation Model.
\bibitem{ref6} Yu Fang, Kanchana Ranasinghe, Le Xue, Honglu Zhou, Juntao Tan, Ran Xu, Shelby Heinecke, Caiming Xiong, Silvio Savarese, Daniel Szafir, Mingyu Ding, Michael S. Ryoo, \& Juan Carlos Niebles (2025). Robotic VLA Benefits from Joint Learning with Motion Image Diffusion.
\bibitem{ref7} Yihao Wang, Pengxiang Ding, Lingxiao Li, Can Cui, Zirui Ge, Xinyang Tong, Wenxuan Song, Han Zhao, Wei Zhao, Pengxu Hou, Siteng Huang, Yifan Tang, Wenhui Wang, Ru Zhang, Jianyi Liu, \& Donglin Wang (2025). VLA-Adapter: An Effective Paradigm for Tiny-Scale Vision-Language-Action Model.
\bibitem{ref8} Haohan Chi, Huan-ang Gao, Ziming Liu, Jianing Liu, Chenyu Liu, Jinwei Li, Kaisen Yang, Yangcheng Yu, Zeda Wang, Wenyi Li, Leichen Wang, Xingtao Hu, Hao Sun, Hang Zhao, \& Hao Zhao (2025). Impromptu VLA: Open Weights and Open Data for Driving Vision-Language-Action Models.
\bibitem{ref9} Shaoqi Dong, Chaoyou Fu, Haihan Gao, Yi-Fan Zhang, Chi Yan, Chu Wu, Xiaoyu Liu, Yunhang Shen, Jing Huo, Deqiang Jiang, Haoyu Cao, Yang Gao, Xing Sun, Ran He, \& Caifeng Shan (2025). VITA-VLA: Efficiently Teaching Vision-Language Models to Act via Action Expert Distillation.
\bibitem{ref10} Yu Fang, Yuchun Feng, Dong Jing, Jiaqi Liu, Yue Yang, Zhenyu Wei, Daniel Szafir, \& Mingyu Ding (2026). When Vision Overrides Language: Evaluating and Mitigating Counterfactual Failures in VLAs.
\bibitem{ref11} Yunchao Ma, Yizhuang Zhou, Yunhuan Yang, Tiancai Wang, \& Haoqiang Fan (2025). Running VLAs at Real-time Speed.
\bibitem{ref12} M E Hemler, C Crouse, Y Takada, \& A Sonnenberg (1988). Multiple very late antigen (VLA) heterodimers on platelets. Evidence for distinct VLA-2, VLA-5 (fibronectin receptor), and VLA-6 structures.. Journal of Biological Chemistry. https://doi.org/10.1016/s0021-9258(18)68549-7
\bibitem{ref13} undefined, , \& undefined (2016). Olympia. Timocratis poena, s. Vla.. Supplementum Epigraphicum Graecum. https://doi.org/10.1163/1874-6772\_seg\_a11\_1191
\bibitem{ref14} undefined, , \& undefined (2011). VLA. SpringerReference. https://doi.org/10.1007/springerreference\_40005
\bibitem{ref15} undefined, , \& undefined (2011). VLA-4. SpringerReference. https://doi.org/10.1007/springerreference\_40006
\bibitem{ref16} undefined, , \& undefined (2011). Very late activation (VLA) molecules. SpringerReference. https://doi.org/10.1007/springerreference\_39994
\bibitem{ref17} M. Ghafourian, F. Ghalambor, \& N. Emad Mosth (2009). An Immunohistochemical Study of Beta1 Integrin Molecules (VLA-4, VLA-5, VLA-6) in All Endometrial Compartments of Fertile and Infertile Women in Ahwaz-Iran. Pakistan Journal of Biological Sciences. https://doi.org/10.3923/pjbs.2009.360.366
\bibitem{ref18} Eileen Fingerman \& Martin E. Hemler (1988). Regulation of proteins in the VLA cell substrate adhesion family: Influence of cell growth conditions on VLA-1, VLA-2, and VLA-3 expression. Experimental Cell Research. https://doi.org/10.1016/0014-4827(88)90031-6
\bibitem{ref19} Hideki Hattori, Shinichi Tagawa, Hirohiko Shibayama, Ryoich Inoue, Shuichi Katagiri, Takashi Machii, \& Teruo Kitani (1996). VLA-5 in the Plasma Cell Line, FR4ds, Acts as a Common Regulator of VLA-4 and VLA-6 in Spreading Induced by Fibronectin and Laminin. Cellular Immunology. https://doi.org/10.1006/cimm.1996.0294
\end{thebibliography}

\end{document}
